\subsection{Package organization}

PyGPLA follows a function-oriented, layered design that mirrors the GPLA workflow,
from preparation of the input data to statistical assessment. The main entry point is
the high-level \texttt{pygpla.api.gpla} function, which coordinates preprocessing,
coupling estimation, singular-value decomposition, and optional significance testing
in one call. This keeps the default workflow compact while preserving access to
lower-level components for advanced analyses, methodological modifications, and
extensions.

Core numerical routines are separated by responsibility. The
\texttt{pygpla.preprocessing} module handles trial concatenation, temporal and unit
selection, spike-count filtering, optional LFP normalization, and the invocation of
optional whitening. The \texttt{pygpla.core} package implements whitening,
coupling-matrix construction, and singular-value decomposition, including
post-processing operations such as phase alignment, spike-vector normalization, and
unwhitening-aware handling of LFP vectors. Statistical procedures are collected in
\texttt{pygpla.stats}, which provides surrogate-based testing using multiple
spike-jittering schemes \citep[e.g., multiple jittering schemes; also, see][]{grun2009data}
and an analytical test based on the Marchenko--Pastur random matrix theory threshold
\citep{safaviUnivariateMultivariateCoupling2021}.

To support reproducibility and method validation, PyGPLA includes dedicated simulation
utilities in \texttt{pygpla.simulations}. These include homogeneous and inhomogeneous
Poisson generators, phase-locked spike simulations, and a transient-coupling simulator
that generates paired spike and LFP signals. The repository additionally provides a
figure-reproduction and visualization script and an automated test suite for validating
individual numerical components and end-to-end behavior. Finally, the configuration
dataclasses in \texttt{pygpla.config} provide structured parameter containers for
preprocessing, whitening, and statistical testing.
